\scanpage{122}
\mspage{3}
\deleted{
$M_v(S')=$ ensemble des $g\in G(S')$ tels que $\deleted{\unsure{g\times_S W}}\quad v_{S'}\in g.W(S')$, i.e. tels que $g^{-1}v_{S'}\in W(S')$. Si $v$ est une section de $W$ et non seulement \inserted{de} $V$, alors $M_v$ contient évidemment $N$.

Dans ces explicitations, on n'a pas utilisé le fait que $G,V,W$ étaient représentables (ni que le site sur lequel on travaille est \deleted{\unsure{e}} défini en termes de préschémas ! ). Mais supposons maintenant que $N$ soit représentable et fidèlement plat et quasi-compact sur $S$, et que $G/N$ soit représentable. Lorsque $S$ est le spectre d'un corps, et que $G$ est de type fini sur $k$, on sait que cette hypothèse est \deleted{\illegible{1}} nécessairement satisfaite (Exp VIII 6 et Exp VI). On voit alors sur le diagramme cartésien (xx), utilisant la théorie de la descente fidèlement plate quasi-compacte et le fait que \inserted{$Z\longrightarrow G\times_S V$} est une immersion fermée, que $X$ est représentable (il est obtenu par descente \deleted{\illegible{1}} du sous-préschéma fermé $Z$ de $G\times_S V$ par le morphisme fidèlement plat et quasi-compact $G\times_S V\longrightarrow G/N\times_S V$). Donc le diagramme (D) est un diagramme de morphismes de \deleted{\unsure{schémas}} préschémas sur $S$.

Nous supposons par la suite que $S$ est le spectre d'un corps $k$, et que $G,V,W$ sont de type fini sur $k$. Soit \deleted{\unsure{noté}} $\uline{\mathfrak{n}}$ l'algèbre de Lie de $N$, donc on a $\dim N\leq\operatorname{rang}\uline{\mathfrak{n}}$, l'égalité étant vraie si et seulement si $N$ est lisse sur $k$ (Exp VI). Soit $a\in W(k)$, et considérons le sous-schéma $M_a$ de $G$ défini plus haut, contenant $N$, et isomorphe à $\phi^{-1}(a)$; nous désignerons par $\uline{\mathfrak{m}_a}$ son espace tangent de Zariski en l'élément neutre $e$ de $N$, de sorte qu'on a
\[
(1)\qquad\uline{\mathfrak{n}\subset\mathfrak{m}_a},\quad\deleted{\unsure{\text{donc}}}\quad\dim N\leq\operatorname{rang}\mathfrak{n}\leq\operatorname{rang}\mathfrak{m}_a.
\]
\uline{Lemme 1.1.} \deleted{\unsure{a)}} Avec les notations précédentes:

a) Considérons les conditions suivantes:
\[
(\mathrm{i})\quad\mathfrak{n}=\mathfrak{m}_a\quad\deleted{\illegible{1}}\quad\text{et }N\text{ est lisse sur }k.
\]
}
\margin{?}
\margin{?}
\margin{121}
